%% notes-tex/010-additive-baselines/sections/2-what-is-predicted.tex

\section{What is being predicted\secstatus{todo}}
\label{sec:label}

\subsection{The label\secstatus{todo}}

Write $f_i$ for the fitness of the single mutant $\Delta i$, $f_{ij}$ for the
double, $f_{ijk}$ for the triple, each relative to wild type. The
\term{digenic interaction score} is the departure of the double from the
multiplicative expectation,
%
\begin{equation}
  \varepsilon_{ij} \;=\; f_{ij} - f_i f_j ,
  \label{eq:epsilon}
\end{equation}
%
and the \term{trigenic interaction score} is the departure of the triple from
the expectation built out of the singles \emph{and} the doubles,
%
\begin{equation}
  \tau_{ijk} \;=\; f_{ijk}
  \;-\; f_i f_j f_k
  \;-\; f_i\,\varepsilon_{jk}
  \;-\; f_j\,\varepsilon_{ik}
  \;-\; f_k\,\varepsilon_{ij} .
  \label{eq:tau}
\end{equation}
%
\external{Kuzmin et al.\ 2018, \texttt{10.1126/science.aao1729}; the published
definition of the released adjusted score column, not re-derived here}

Two properties of Eq.~\ref{eq:tau} matter for everything that follows. It is a
\emph{residual}: the additive-on-the-log-scale, multiplicative-on-the-linear-scale
part of fitness is subtracted before the number is reported. And it is a
\emph{third-order} residual: the pairwise terms are subtracted too, so $\tau$ is
by construction the part of the triple mutant phenotype that neither single nor
double mutant measurements predict.

The 010 loader consumes the column \texttt{Adjusted genetic interaction score
(epsilon or tau)} from the Kuzmin releases, which for a trigenic record is
$\tau$, and stores it under the label key \texttt{gene\_interaction}.

\subsection{The build\secstatus{todo}}

%% SOURCE: read from the pinned build by
%% experiments/010-kuzmin-tmi/scripts/additive_baseline_gene_interaction.py
\begin{table}[htbp]
\centering
\begin{tabular}{ll}
\hline
Records & 376{,}732 \\
Perturbation count & 3 for every record \\
Distinct perturbed genes & 4{,}352 \\
Sources & \texttt{TmiKuzmin2018Dataset}, \texttt{TmiKuzmin2020Dataset} \\
Label & \texttt{gene\_interaction} ($\tau$) \\
Split, seed 42 & train 301{,}386 / val 37{,}673 / test 37{,}673 \\
Train label SD & 0.063186 \\
Val label SD & 0.062914 \\
Test label SD & 0.064187 \\
\hline
\end{tabular}
\caption{The 010 build. Every baseline and the transformer use these records,
this split and this label. Read from
\texttt{001-small-build/processed/label\_df.parquet},
\texttt{is\_any\_perturbed\_gene\_index.json} and
\texttt{data\_module\_cache/index\_seed\_42.json}.}
\label{tab:build}
\end{table}

The build is triples only, so there are no singles or doubles in it. The
MULTI-evolve experiment of training on low order and extrapolating to high order
is therefore not available here, and the question is the narrower one: given a
training set of triples, how much of the held-out $\tau$ is recoverable without
representing interaction.

\subsection{The screen design, which is the confounder\secstatus{todo}}

Kuzmin's trigenic screen crosses a \term{query} double mutant against an array of
single mutants, so a triple is a query pair plus one array gene, and the same
query pair recurs across hundreds or thousands of triples. Counting unordered
gene pairs in the training split finds 593{,}978 distinct pairs, of which only
420 occur five or more times, the most frequent occurring 3{,}308 times. That
concentration is the query-pair structure.

The 010 split is random over records, so a query pair that appears in training
also appears in validation and test. Any model that can represent a per-pair
offset can absorb whatever is common to a query batch, and none of that is
trigenic biology. Baseline B4 in Section~\ref{sec:models} is built to measure
exactly this quantity.
